\section{Data model OLAP}\label{sec:data-model}
\subsection{Thiết kế kho dữ liệu}
\subsubsection*{Bảng fact và mô hình vật lý}
Kho dữ liệu có sáu bảng fact thuộc cả ba loại của Kimball. Hai bảng giao dịch ở mức chi tiết nhất là \code{fact\_\allowbreak{}order\_\allowbreak{}item} (một dòng cho mỗi sản phẩm trong đơn, 112.650 dòng khi nạp đủ) và \code{fact\_\allowbreak{}payment} (một dòng cho mỗi lần thanh toán); \code{fact\_\allowbreak{}review} ghi điểm đánh giá, kèm cờ cho biết đơn có giao trễ hay không. Bảng \code{fact\_\allowbreak{}order\_\allowbreak{}fulfillment} là snapshot tích luỹ: mỗi đơn một dòng và được cập nhật khi đơn đi qua từng mốc (mua, duyệt, bàn giao vận chuyển, giao khách, ngày dự kiến), từ đó tính thời gian từng chặng, số ngày trễ và số ngày Olist hứa giao. Quy trình thu hút seller dùng snapshot tích luỹ thứ hai \code{fact\_\allowbreak{}lead\_\allowbreak{}funnel} (mỗi khách tiềm năng một dòng, đi từ liên hệ đầu, chốt deal đến đơn đầu tiên), còn \code{fact\_\allowbreak{}seller\_\allowbreak{}monthly} là snapshot định kỳ, mỗi seller mỗi tháng một dòng, lưu dày đặc để đo được cả những tháng seller không bán. Mọi khoá ngoại đều không rỗng; nếu không tra được chiều thì dùng dòng Unknown và ghi vào nhật ký ETL. Khoá khách hàng luôn trỏ tới phiên bản địa chỉ có hiệu lực vào ngày mua. Các khoảng thời gian âm do lỗi nguồn được đặt rỗng và đếm lại, không xoá đơn. Sau khi nạp, tổng giá bán ở fact khớp từng đồng với dữ liệu nguồn (13.591.643,70 BRL khi nạp đủ), và 112 phép kiểm thử về số dòng, tổng tiền, grain, khoá ngoại và logic nghiệp vụ đều đạt. Ở mức vật lý, mỗi bảng fact có khoá chính theo grain, khoá ngoại sang các chiều và chỉ mục B-tree trên mọi khoá ngoại.

\begin{table}[htbp]
\centering\small
\caption{Các bảng fact: loại, grain và số dòng}
\label{tab:facts}
\begin{tabular}{lllr}
\toprule
\textbf{Fact} & \textbf{Loại} & \textbf{Grain} & \textbf{Số dòng} \\
\midrule
fact\_order\_item & Transaction & 1 dòng hàng của 1 đơn & 112.650 \\
fact\_order\_fulfillment & Accumulating snapshot & 1 đơn (5 mốc thời gian) & 99.441 \\
fact\_payment & Transaction & 1 lần thanh toán & 103.886 \\
fact\_review & Transaction & 1 review của 1 đơn & 99.224 \\
fact\_lead\_funnel & Accumulating snapshot & 1 MQL & 8.000 \\
fact\_seller\_monthly & Periodic snapshot (dày đặc) & 1 seller × 1 tháng & 36.583 \\
\bottomrule
\end{tabular}
\end{table}

\hinh{3-3-4_erd_dw.png}{Mô hình dữ liệu của kho (data model; bản nháp tự động)}{fig:erd-dw}{1.0}
\hinh{3-3-4_constellation.png}{Mô hình constellation theo data mart: fact (cam), dimension (xanh) (bản nháp tự động)}{fig:constellation}{0.9}
\todo{Yêu cầu của môn: \textbf{data model logic và vật lý} của staging / data mart / kho, dùng màu phân biệt fact (cam) và dimension (xanh), nhóm các vùng liên quan bằng khung. Vẽ lại bằng draw.io từ docs/diagrams/data-model-theo-mart.drawio. Thêm: hệ thống chiều khái niệm (phân cấp từng dimension), thống kê số bản ghi từng dimension (đã có Bảng \ref{tab:dims}), phân nhóm các dimension.}
\todo{Thêm bus matrix (docs/DW\_DESIGN.md mục 4) và sơ đồ phân cấp: ngày, địa lý, sản phẩm, seller, khách (SCD2).}

\subsection{Data mart và các thao tác OLAP}
Trên kho dữ liệu hình sao, nhóm dựng bốn data mart theo chủ đề: bán hàng, xử lý đơn, trải nghiệm khách và thu hút seller. Mỗi mart là một lớp view trong schema riêng, tổng hợp sẵn các độ đo mà dashboard cần (GMV theo ngày, danh mục, bang; thời gian từng chặng và tỷ lệ giao trễ; điểm đánh giá theo mức trễ, cơ cấu thanh toán, RFM, cohort; funnel từ khách tiềm năng đến seller có đơn và scorecard seller). Các view có chi phí tính toán lớn được cài dưới dạng materialized view và được làm mới trong mỗi lần chạy pipeline, nên mọi truy vấn trên mart đều trả lời dưới 0,2 giây. Các thao tác OLAP được minh hoạ bằng SQL theo mô hình ROLAP: roll-up bằng ROLLUP, GROUPING SETS và CUBE; drill-down bằng cách thêm cấp thấp hơn của phân cấp vào GROUP BY; slice bằng cách cố định một chiều (năm 2017); dice bằng cách giới hạn đồng thời nhiều chiều (vùng Sudeste, nhóm ngành Home \& Furniture, quý 4); pivot bằng tổng có điều kiện FILTER. Ví dụ, lát cắt năm 2017 cho thấy vùng Sudeste chiếm 64,3\% GMV, và drill-down tháng 11/2017 chỉ ra ngày Black Friday 24/11 có GMV cao nhất tháng (152.653,74 BRL). Nhóm chọn ROLAP thay vì MOLAP vì hệ thống dùng PostgreSQL, dữ liệu vài trăm nghìn dòng fact, kho luôn cập nhật sau mỗi lần nạp tăng dần mà không cần xử lý lại khối đa chiều, còn Power BI đảm nhiệm lớp ngữ nghĩa và tương tác kéo thả cho người dùng cuối. Lớp ngữ nghĩa đó gồm 32 quan hệ, các phân cấp thời gian, địa lý, sản phẩm, seller và thư viện measure DAX cho các KPI chính, được đối chiếu với giá trị tính bằng SQL.

{\small\begin{longtable}{llllp{6.3cm}}
\caption{Các data mart: view, loại, số dòng và mô tả}\label{tab:marts}\\
\toprule
\textbf{Mart} & \textbf{View} & \textbf{Loại} & \textbf{Dòng} & \textbf{Mô tả} \\
\midrule
\endfirsthead
\toprule
\textbf{Mart} & \textbf{View} & \textbf{Loại} & \textbf{Dòng} & \textbf{Mô tả} \\
\midrule
\endhead
\bottomrule
\endfoot
mart\_acquisition & v\_funnel\_by\_origin\_month & view & 113 & Funnel theo tháng liên hệ đầu x origin: MQL, đã chốt, đã có đơn, tỷ lệ, số ngày TB, GMV 90 ngày đầu \\
mart\_acquisition & v\_seller\_by\_channel & materialized & 7 & Theo kênh thu hút seller: số seller, số seller có đơn, GMV, GMV / seller, trễ, review; cờ mẫu nhỏ (< 30 seller có đơn) \\
mart\_acquisition & v\_seller\_scorecard & materialized & 3.095 & 1 dòng / seller có đơn: GMV, đơn, trễ, hủy, review, hạng và tỷ trọng luỹ kế (Pareto), cờ mẫu nhỏ (< 10 đơn) \\
mart\_cx & v\_customer\_cohort & materialized & 220 & Cohort theo tháng mua đầu tiên: số khách còn mua lại sau 0,1,2... tháng và tỷ lệ quay lại \\
mart\_cx & v\_customer\_rfm & materialized & 94.983 & 1 dòng / khách thật (customer\_unique\_id): recency, frequency, monetary, điểm R-F-M 1-5 và phân khúc; đơn không thuộc nhóm Failed \\
mart\_cx & v\_payment\_mix\_month & view & 90 & Cơ cấu thanh toán theo tháng mua x hình thức: số lần, giá trị, trả góp TB, tỷ trọng giá trị \\
mart\_cx & v\_review\_vs\_delay & view & 105 & Điểm review và \% tiêu cực theo nhóm trễ giao (giao sớm ... trễ trên 7 ngày) x tháng mua; chỉ đơn đã giao \\
mart\_fulfillment & v\_fulfillment\_order & view & 99.441 & 1 dòng / đơn: mốc ngày, thời gian từng chặng, số ngày hứa và trễ, nhóm trễ (delay\_bucket), bang và vùng của khách \\
mart\_fulfillment & v\_fulfillment\_state\_month & view & 565 & Theo tháng mua x bang: số đơn, tỷ lệ đúng hạn / hủy, thời gian trung bình từng chặng (chỉ đơn đã giao) \\
mart\_fulfillment & v\_shipping\_distance & view & 112.650 & Khoảng cách seller -> khách (km, haversine) cho từng dòng hàng, nhóm 0-100 / 100-500 / 500-1000 / trên 1000 km; distance\_km NULL nếu thiếu toạ độ (has\_distance = false) \\
mart\_fulfillment & v\_status\_funnel & view & 125 & Số đơn đi qua từng chặng (đặt, duyệt, bàn giao, giao, hủy) theo tháng mua, dạng dài \\
mart\_sales & v\_category\_pareto & view & 74 & Pareto 80/20 của GMV theo danh mục: xếp hạng, tỷ trọng, tỷ trọng luỹ kế, hạng A/B/C \\
mart\_sales & v\_geo\_state & view & 27 & Dữ liệu bản đồ theo bang: toạ độ tâm (trung bình các zip có toạ độ), số khách, đơn, GMV, freight ratio, số seller, lead time, tỷ lệ trễ \\
mart\_sales & v\_geo\_zip & view & 14.889 & Dữ liệu bản đồ theo zip có toạ độ và có khách hoặc seller: toạ độ, thành phố, số đơn và GMV phía khách, số seller \\
mart\_sales & v\_sales\_category\_month & view & 1.282 & GMV, đơn, dòng hàng theo tháng x nhóm ngành hàng x danh mục (số đơn không cộng được giữa các danh mục) \\
mart\_sales & v\_sales\_daily & view & 614 & GMV, số đơn, số dòng hàng, phí vận chuyển, AOV theo ngày mua (chỉ đơn không thuộc nhóm Failed) \\
mart\_sales & v\_sales\_state\_month & view & 558 & GMV, đơn, phí vận chuyển, tỷ lệ freight/price theo tháng x bang của khách \\
\end{longtable}}

\paragraph{Roll-up và pivot trên dữ liệu thật.} Bảng \ref{tab:olap-rollup} là kết quả roll-up bằng \code{ROLLUP} (SQL ở Phụ lục B); Bảng \ref{tab:olap-pivot} là pivot GMV theo vùng và năm bằng tổng có điều kiện \code{FILTER}.
\begin{table}[htbp]
\centering\small
\caption{Roll-up bằng ROLLUP: GMV theo quý, năm và tổng}
\label{tab:olap-rollup}
\begin{tabular}{llrr}
\toprule
\textbf{Năm} & \textbf{Quý} & \textbf{Số đơn} & \textbf{GMV (BRL)} \\
\midrule
2017 & 2017-Q1 & 5.122 & 733.398,94 \\
2017 & 2017-Q2 & 9.222 & 1.286.918,78 \\
2017 & 2017-Q3 & 12.445 & 1.681.949,00 \\
2017 & 2017-Q4 & 17.586 & 2.406.225,55 \\
2017 & (cả năm) & 44.375 & 6.108.492,27 \\
2018 & 2018-Q1 & 20.979 & 2.764.402,78 \\
2018 & 2018-Q2 & 19.897 & 2.849.730,26 \\
2018 & 2018-Q3 & 12.655 & 1.727.049,37 \\
2018 & (cả năm) & 53.531 & 7.341.182,41 \\
Tổng & (cả năm) & 97.906 & 13.449.674,68 \\
\bottomrule
\end{tabular}
\end{table}

\begin{table}[htbp]
\centering\small
\caption{Pivot bằng FILTER: GMV (BRL) theo vùng × năm}
\label{tab:olap-pivot}
\begin{tabular}{lrrrr}
\toprule
\textbf{Vùng} & \textbf{2016} & \textbf{2017} & \textbf{2018} & \textbf{Tổng} \\
\midrule
Sudeste & 27.701 & 3.928.343 & 4.866.293 & 8.822.336 \\
Sul & 7.456 & 883.020 & 1.047.636 & 1.938.111 \\
Nordeste & 6.013 & 727.733 & 804.568 & 1.538.314 \\
Centro-Oeste & 2.356 & 401.380 & 458.090 & 861.825 \\
Norte & 1.200 & 168.018 & 164.596 & 333.814 \\
\bottomrule
\end{tabular}
\end{table}

\todo{Chọn thêm các truy vấn drill-down, slice, dice từ handoff/olap-examples.md; viết đoạn so sánh ROLAP và MOLAP trong bối cảnh bài toán (liên hệ mục 1.6).}

\subsection{Hiệu năng truy vấn}
Nhóm đo thời gian của bốn truy vấn phân tích trên dữ liệu đầy đủ (112.650 dòng hàng) bằng EXPLAIN ANALYZE, lấy trung bình năm lần. Trên các truy vấn quét toàn bộ dữ liệu, lược đồ hình sao chỉ nhanh hơn lược đồ chuẩn hoá của hệ thống nguồn từ 1,0 đến 1,6 lần vì dữ liệu còn nhỏ; lợi ích chính của mô hình chiều là truy vấn đơn giản và nhất quán. Chỉ mục trên khoá ngoại không thay đổi các truy vấn quét toàn bộ nhưng làm truy vấn chọn lọc theo một seller nhanh gấp khoảng chín lần. Đáng kể nhất là materialized view: đọc kết quả tổng hợp sẵn mất dưới một mili giây so với 150–750 mili giây khi view thường phải tính lại, với chi phí là phải làm mới mỗi lần nạp dữ liệu (khoảng hai giây). Fact snapshot định kỳ theo seller-tháng cũng làm scorecard seller nhanh gấp hai lần so với tính từ dòng hàng.

\begin{table}[htbp]
\centering\small
\caption{Thời gian thực thi trung bình của 4 truy vấn (mili giây, trung bình 5 lần EXPLAIN ANALYZE)}
\label{tab:benchmark}
\begin{tabular}{p{4.6cm}rrrrr}
\toprule
\textbf{Truy vấn} & \textbf{oltp (chuẩn hoá)} & \textbf{dw (star, có index FK)} & \textbf{dw không index FK} & \textbf{View thường} & \textbf{Materialized view} \\
\midrule
Q1 GMV theo bang × tháng & 657.4 & 506.1 & 494.0 & 770.4 & 0.056 \\
Q2 Lead time theo bang & 85.7 & 53.8 & 51.5 & 165.8 & 0.042 \\
Q3 Scorecard seller & 375.1 & 386.8 & 391.3 & 196.3 & 0.193 \\
Q4 GMV của một seller (tra cứu chọn lọc) & 6.3 & 0.8 & 7.3 & -- & -- \\
\bottomrule
\end{tabular}
\end{table}

\hinh{3-5-1_benchmark.png}{Thời gian thực thi trung bình: oltp so với dw, có và không có index khóa ngoại, view thường so với materialized view (thang logarit)}{fig:benchmark}{1.0}
\paragraph{Nhận xét.} Trên dữ liệu khoảng 100 nghìn đơn, lược đồ hình sao chỉ nhanh hơn lược đồ chuẩn hóa của hệ thống nguồn từ 1,0 đến 1,6 lần ở các truy vấn quét toàn bộ, và scorecard seller không nhanh hơn; lợi ích chính của mô hình chiều là \textbf{truy vấn đơn giản và số liệu nhất quán} (một định nghĩa cho mỗi KPI). Tốc độ đến chủ yếu từ hai kỹ thuật: \textbf{materialized view} (đọc dưới 1 ms thay vì 150--750 ms) và \textbf{chỉ mục} trên khóa ngoại (truy vấn chọn lọc một seller nhanh gấp khoảng 9 lần). Cần lưu ý đo trên một máy, một người dùng, dữ liệu còn nhỏ.
\todo{Nhóm xác nhận số đo trên máy của mình nếu chạy lại (python scripts/benchmark.py); thêm giải thích kế hoạch thực thi (EXPLAIN) cho một truy vấn nếu muốn.}
